\documentclass[a4paper,11pt]{article}

\usepackage{fontspec}
\setmainfont{Carlito}
\usepackage[margin=0.6in, top=0.5in, bottom=0.5in]{geometry}
\usepackage{titlesec}
\usepackage{enumitem}
\usepackage{xcolor}
\usepackage[colorlinks=true]{hyperref}
\usepackage{microtype}
\usepackage{ragged2e}
\usepackage{needspace}

% ---- palette: one accent, one ink, one muted grey ----
\definecolor{navy}{RGB}{12, 83, 166}
\definecolor{ink}{RGB}{33, 37, 41}
\definecolor{muted}{RGB}{96, 102, 110}
\hypersetup{urlcolor=navy, linkcolor=navy}
\color{ink}

% ---- ragged right, no hyphenation: even word spacing, no split words ----
\hyphenpenalty=10000 \exhyphenpenalty=10000 \emergencystretch=3em
\setlength{\parindent}{0pt}
\pagestyle{empty}
\linespread{1.03}
\AtBeginDocument{\RaggedRight}

% ---- type scale (ratio ~1.2): body 11 -> headline 12 -> section 13 -> name 22 ----
\titleformat{\section}
  {\color{navy}\fontsize{13}{15}\selectfont\bfseries\addfontfeature{LetterSpace=6}}
  {}{0em}{\MakeUppercase}[{\color{navy}\vspace{1pt}\titlerule[0.8pt]}]
\titlespacing*{\section}{0pt}{9pt}{5pt}

\setlist[itemize]{label={\color{navy}\small\textbullet}, leftmargin=1.1em, labelsep=0.45em,
  itemsep=1.5pt, parsep=0pt, topsep=2pt, partopsep=0pt}

% ---- building blocks: bold = organisation/project, italic = role/stack, grey = dates/places ----
\newcommand{\role}[4]{% org, location, dates, title
  \textbf{#1}{\color{muted}\,\textperiodcentered\,\textit{#2}}\hfill{\color{muted}#3}\\
  \textit{#4}\par}
\newcommand{\roleNoLoc}[3]{\textbf{#1}\hfill{\color{muted}#2}\\ \textit{#3}\par}
\newcommand{\project}[3]{% name, stack, right-hand tag
  \textbf{#1}\hfill#3\\ {\color{muted}\small\textit{#2}}\par}
\newcommand{\gh}[1]{\href{#1}{\small GitHub}}
\newcommand{\badge}[1]{{\color{navy}\small\textbf{#1}}}
\newcommand{\skillgroup}[2]{{\color{navy}\textbf{#1}}\\ #2\par\vspace{3pt}}
\newcommand{\gap}{\vspace{5pt}}
\newcommand{\num}[1]{\textbf{#1}}   % result numbers only


\begin{document}

% ================= HEADER =================
\begin{center}
  {\fontsize{22}{26}\selectfont\bfseries\color{navy} Thiruvel Andagurunathan Pandian}\\[4pt]
  {\fontsize{12}{14}\selectfont\color{navy} Machine Learning Engineer}\\[4pt]
  {\small
  Bristol, UK\nobreak\ \,\textperiodcentered\, \href{mailto:apthiruvel@gmail.com}{apthiruvel@gmail.com}\nobreak\ \,\textperiodcentered\, \href{tel:+447741016954}{+44 7741 016954}\nobreak\ \,\textperiodcentered\, \href{https://www.linkedin.com/in/thiruvel-a-p/}{linkedin.com/in/thiruvel-a-p}\nobreak\ \,\textperiodcentered\, \href{https://github.com/Thiruvel-AP}{github.com/Thiruvel-AP}\nobreak\ \,\textperiodcentered\, \href{https://thiruvel-ap.github.io/PersonalPortfolio/}{Portfolio}}\\[3pt]
  {\small\itshape\color{muted} Right to Work in UK\nobreak\ \,\textperiodcentered\, Available to join Immediately}

\end{center}

% ================= PROFILE =================
\needspace{5\baselineskip}\section{Profile}
At 180 Degrees Consulting, led the architecture for a product startup client\textquotesingle{}s hierarchical agent workflows (PostgreSQL for persistent storage and retrieval across agent sessions), demoed it live (five star rating from senior leadership) and contributed to the client\textquotesingle{}s production TypeScript/Node.js MCP server (OpenAI based query generation), now in live use. Uses Claude Code, Cursor and Codex daily. Owns data workflows end to end in Python and SQL, from merging structured, API and web scraped data through feature creation, tuning, validation and deployment (Docker, GitHub Actions CI/CD, MLflow, AWS Lambda). Shipped production fintech features integrating Twelve Data market feeds and Plaid bank linking. 3rd place, Lloyd\textquotesingle{}s Banking Group Datathon.

% ================= SKILLS =================
\needspace{5\baselineskip}\section{Skills}
\skillgroup{Agentic AI \& LLM Applications}{Tool calling and failure recovery (LangChain, LangGraph, Google ADK, ReAct agent loops; circuit breakers, fallbacks, deterministic guardrails), multi provider LLMs (OpenAI, Anthropic Claude on AWS Bedrock, Google Gemini via LiteLLM, open weight models via Ollama), LLM evaluation (three tier tests, hallucination rate analysis), agent memory (persistent storage and retrieval across sessions), RAG pipelines (retrieval, generation, embeddings and the ChromaDB vector database), MCP servers (FastMCP, TypeScript SDK), prompt engineering, AI assisted development (Claude Code as primary tool, Cursor, OpenAI Codex)}
\skillgroup{Software Engineering \& Integrations}{TypeScript/Node.js, React, Python, JavaScript, FastAPI, Pydantic, RESTful APIs and third party API integrations, GraphQL, PostgreSQL, SQL, MongoDB and DynamoDB (NoSQL), graph data (networkx), pytest, Git (pull requests, code review), debugging and root cause analysis}
\skillgroup{Cloud \& Deployment}{AWS (Bedrock, Lambda, ECR, S3, SQS, EC2, IAM, Cognito), GCP (Cloud Run), Docker, Kubernetes (AWS EKS), GitHub Actions CI/CD, serverless deployment}

% ================= EXPERIENCE =================
\needspace{5\baselineskip}\section{Work Experience}
\needspace{4\baselineskip}
\role{AstraZeneca}{Bristol, UK (Remote)}{May 2026 – Sep 2026}{Data Scientist}
\begin{itemize}
  \item Needed an LLM a research team could trust before it shipped: built a tool calling LLM agent (LangChain, AWS Bedrock, Claude Haiku 4.5) for a multiomics cancer research platform, with three tools for gene alias resolution, scoring methodology explanation and dataset provenance.
  \item Designed a deterministic hallucination guard architecture and fail safe data resolution via live Ensembl/HGNC APIs with a local \num{19,213 gene parquet fallback} and circuit breakers.
  \item Diagnosed where it struggled with a three tier, \textasciitilde{}65 test evaluation harness (tool level → API → live) and hallucination rate analysis, debugging failures layer by layer, and used the same harness to validate the Groq to Bedrock migration after a rate limit evaluation.
  \item Built the Python data pipeline: reconciled \num{14 heterogeneous source datasets} through identifier crosswalks so the agent used trusted data, then deployed it as a Docker image on AWS Lambda (ECR, least privilege IAM) and documented the architecture and evaluation.
\end{itemize}
\gap
\needspace{4\baselineskip}
\role{180 Degrees Consulting Bristol}{Bristol, UK (Remote)}{Sep 2025 – Dec 2025}{Data Consultant (AI Engineering \& AI System)}
\begin{itemize}
  \item A product startup client needed an agentic AI prototype turned into software their team could use: led the architecture for a five person cross functional team, designed hierarchical agent workflows (MCP, Exograph GraphQL, PostgreSQL for persistent storage and retrieval across agent sessions) and built the Python prototype (FastMCP, Ollama, GitHub Actions CI), then ran a live demo with client stakeholders \num{(five star rating from senior leadership)} within the three month engagement.
  \item Proposed the request flow: a user\textquotesingle{}s text goes through the API and a customised prompt to the MCP server, where one LLM instance writes GraphQL queries against the client\textquotesingle{}s Convex database and a second instance runs the statistical analysis before the result returns to the user.
  \item Carried it into production with the client\textquotesingle{}s engineering team: contributed to their production TypeScript/Node.js MCP server (OpenAI based query generation, Convex), now in live use, and reviewed pull requests in the client\textquotesingle{}s production repository alongside their CTO.
\end{itemize}
\gap
\needspace{4\baselineskip}
\role{Freelance}{Chennai, India}{Jan 2025 – Aug 2025}{AI/ML Engineer}
\begin{itemize}
  \item Built scikit-learn pipelines and LSTM/BiRNN/CNN models in TensorFlow/Keras, deployed as FastAPI services in Docker.
  \item Built NLP and agentic RAG workflows with Hugging Face Transformers, Ollama, LangChain and LangGraph.
  \item Designed reproducible scikit-learn pipelines with feature engineering (temporal features, interaction terms, target encoding), model selection, cross validation and hyperparameter tuning across classification and regression tasks.
\end{itemize}
\gap
\needspace{4\baselineskip}
\role{Avasoft}{Chennai, India}{Jan 2023 – Sep 2024}{Software Development Engineer (Jun 2023 – Sep 2024); Software Development Engineer Trainee (Jan 2023 – May 2023)}
\begin{itemize}
  \item Implemented offline first data persistence (Isar) with background synchronisation in the Flutter FinTech app, keeping financial records available in low connectivity conditions.
  \item Shipped features in a production fintech (financial services) application as part of a team: integrated Twelve Data (real time market feeds) and Plaid (bank linking, transaction imports) via REST APIs, with SSL pinning and AWS Cognito SSO.
  \item As a trainee, modelled user profile, transaction history and event booking data in Core Data schemas with CRUD and session state management, and automated a manual data collection process with a JavaScript web scraping pipeline that surfaced market trend signals used for pricing decisions.
\end{itemize}

% ================= PROJECTS =================
\needspace{5\baselineskip}\section{Projects}
\needspace{4\baselineskip}
\setbox0=\hbox{\textbf{AgenticFriend: real time voice assistant}\hspace{1em}\url{https://github.com/Thiruvel-AP/adk_multi-agent}}%
\setbox2=\hbox{{\small\textit{FastAPI, Google ADK, React, faster-whisper, PyTorch, Docker/CUDA}}\hspace{1em}{\small\url{https://github.com/Thiruvel-AP/adk_multi-agent}}}%
\ifdim\wd0>\linewidth
\ifdim\wd2>\linewidth
\textbf{AgenticFriend: real time voice assistant}\\ {\color{muted}\small\textit{FastAPI, Google ADK, React, faster-whisper, PyTorch, Docker/CUDA}}\\ \hfill{\small\url{https://github.com/Thiruvel-AP/adk_multi-agent}}\par
\else
\textbf{AgenticFriend: real time voice assistant}\\ {\color{muted}\small\textit{FastAPI, Google ADK, React, faster-whisper, PyTorch, Docker/CUDA}}\hfill{\small\url{https://github.com/Thiruvel-AP/adk_multi-agent}}\par
\fi
\else
\project{AgenticFriend: real time voice assistant}{FastAPI, Google ADK, React, faster-whisper, PyTorch, Docker/CUDA}{\url{https://github.com/Thiruvel-AP/adk_multi-agent}}
\fi
\begin{itemize}
  \item Text to speech on CUDA in float16 with startup warm up, packaged in an nvidia/cuda Docker image.
\end{itemize}
\gap
\needspace{4\baselineskip}
\setbox0=\hbox{\textbf{CoverFit: insurance benchmarking API (Bristol × UWE AI Hackathon)}\hspace{1em}\url{https://github.com/Thiruvel-AP/AIHackathon-Phoenix}}%
\setbox2=\hbox{{\small\textit{FastAPI, pandas, LiteLLM, Pydantic, pytest}}\hspace{1em}{\small\url{https://github.com/Thiruvel-AP/AIHackathon-Phoenix}}}%
\ifdim\wd0>\linewidth
\ifdim\wd2>\linewidth
\textbf{CoverFit: insurance benchmarking API (Bristol × UWE AI Hackathon)}\\ {\color{muted}\small\textit{FastAPI, pandas, LiteLLM, Pydantic, pytest}}\\ \hfill{\small\url{https://github.com/Thiruvel-AP/AIHackathon-Phoenix}}\par
\else
\textbf{CoverFit: insurance benchmarking API (Bristol × UWE AI Hackathon)}\\ {\color{muted}\small\textit{FastAPI, pandas, LiteLLM, Pydantic, pytest}}\hfill{\small\url{https://github.com/Thiruvel-AP/AIHackathon-Phoenix}}\par
\fi
\else
\project{CoverFit: insurance benchmarking API (Bristol × UWE AI Hackathon)}{FastAPI, pandas, LiteLLM, Pydantic, pytest}{\url{https://github.com/Thiruvel-AP/AIHackathon-Phoenix}}
\fi
\begin{itemize}
  \item An insurance sponsor needed a fast, low cost, explainable benchmark (hackathon brief): paired a deterministic matching engine (FastAPI, pandas, Pydantic) with a single LLM call (Gemini via LiteLLM) for the rationale, covered by pytest tests; responses arrive in \num{under 3 seconds at \textasciitilde{}£0.003 per query}.
\end{itemize}

% ================= EDUCATION =================
\needspace{5\baselineskip}\section{Education}
\needspace{3\baselineskip}
\roleNoLoc{University of Bristol}{Sep 2025 – Sep 2026}{MSc Data Science, Distinction \textit{(Expected)}}
{\color{muted}\small Large scale Data Engineering, Statistical Computing \& Empirical Methods, AI \& Text Analytics, Visual Analytics, Data Science Methods \& Practices, Technology, Innovation, Business \& Society}\par
\gap
\needspace{3\baselineskip}
\roleNoLoc{St. Joseph\textquotesingle{}s College of Engineering}{Aug 2019 – Apr 2023}{B.E. Electrical and Electronics Engineering, 91\%, First Class with Distinction}

% ================= CERTIFICATIONS & ACHIEVEMENTS =================
\needspace{5\baselineskip}\section{Certifications \& Achievements}
\begin{itemize}
  \item Model Context Protocol: Advanced Topics, Anthropic, \textit{Jun 2026}
  \item Introduction to Model Context Protocol, Anthropic, \textit{Jun 2026}
  \item AWS Academy Cloud Foundations, \textit{Nov 2025}
  \item 180 Degrees Consulting Data Science Consultant Certificate, \textit{Dec 2025}
  \item Bristol PLUS Award, \textit{Mar 2026}
\end{itemize}

\end{document}